\chapter{Background and Related Work}
\label{ch2}

\graphicspath{{Chapter2/}}

\section{Software Modeling and Model Driven Engineering}
\label{sec:back_mde}

Model-Driven Engineering (MDE) is a software development paradigm that places model artifacts as the key components of the development process instead of considering them as simple documentation assistants \cite{schmidt2006mde}. In MDE, a metamodel is used to define a modeling language and its abstract syntax, intended as the set of concepts that the language can express, and the constraints that models must satisfy; concrete models that respect the specifications of a metamodel are used to represent specific systems or domains. Model transformation, instead, is a process that allows the automatic translation of models between different formats and representations, enabling the execution of activities such as code generation, model refinement, or migration to different notations. The dominant notation in the MDE ecosystem is the Unified Modeling Language (UML), a graphical notation commonly adopted throughout different steps of the software engineering process by offering diagram families that cover both structural and behavioral aspects of a system, supporting activities such as requirements definition, system and architecture design, use case specification, and activity sequencing. 

Among the various types of diagrams covered by UML, class diagrams are one of the most important and commonly used notations due to their effectiveness for requirements engineering and object-oriented development: through class diagrams it is possible to represent the concepts that compose a system, its design expressed as the combination of hardware and software subsystems, as well as the software classes to implement in code.

This is achieved with the use of \textit{Classes}, rectangle-shaped elements used to represent sets of objects with common properties and relationships, with each specific instance of a class being called an \textit{Object}; properties are called \textit{Attributes} and are characterized by a unique (for the class) name and a type, with each object having a specific value for each defined attribute; the final element expresses relationships between classes through \textit{Associations}, lines that connect two classes and specify a value named \textit{Multiplicity} for both classes: this number states the minimum and maximum number of objects that can be connected in a relationship.

Beyond class diagrams, other key notations defined by UML support the different phases of software development. One of these notations is the concept of use case diagrams, a notation used in requirement engineering processes to capture the functional capabilities of a system from the perspective of the external environment, that is, its intended users and other systems that interact with it. A use case diagram identifies the \textit{actors}, external entities that interact with a given system, the \textit{use cases}, discrete functional units that implement the system's capabilities, and three type of relationships: \textit{associations} that connect actors to the use cases they participate in, either as the executing actor or as a supporter, \textit{inclusion} relationships where one use case unconditionally invokes another as part of its execution flow, and \textit{extension} relationships that identify scenarios where one use case may continue the functionalities of another when certain conditions are verified \cite{fowler2003uml}. The notation's intended use is communication between technical and non-technical stakeholders during the early phases of requirements definition, making it more abstract and purposely less detailed in terms of variability compared to class diagrams, which serve multiple purposes in MDE activities.

Despite being the most commonly used notation for object-oriented design and documentation, UML tends to be underused in both industrial context and academic education, with several studies documenting causes for this gap. A survey asking practitioners and students for perceived difficulties identified limitations such as notation complexity, the necessity of putting effort to keep consistency between diagrams and code, and absence of adequate tools for supporting interactive learning \cite{siau2006identifying}; another investigation on university programming courses conceptualized the issues for UML class diagrams, finding systematic misunderstandings related to object identity, inheritance understanding, and difficulties in representing association semantics that persisted after the learning phase \cite{boustedt2012students}. An observation focused more on the practitioner side found that, during agile development processes, there is a tendency to substitute formal UML representations with informal whiteboards and other, less detailed notations, due to a preference for speed and flexibility against completeness and precision \cite{stettina2012documentation}; in addition to the aforementioned limitations, practitioners and instructors also report barriers such as excessive tool complexity and an inability to perceive enough return on investment to justify adopting UML instead of less formal approaches \cite{oliveirajr2025uml}. The same survey also pointed out misalignment between which diagram types are prominent in educational syllabi, pointing out the need for educational approaches contextualized in real industrial practices. 

The low usage of the UML notation has relevant implications on educational aspects: Bogdanova \cite{bogdanova2019towards} showed that a substantial amount of practice is necessary to reach a basic level of proficiency in conceptual modeling, and identified recurring error patterns among learners over the years; the conclusion behind the study is that sustained and personalized feedback is necessary to help students reach a sufficient level of competence: however, providing such assistance is a complex endeavor in large-scale university courses. A review of tools by Huber and Hagel \cite{huber2022tool} for assisting UML education identified automated and interactive approaches as potentially one of the most effective strategies for improving the learning process, while also pointing out the importance of the quality and immediacy of the feedback on making educational tools actually effective. As a consequence, the definition of feedback mechanisms tailored to students that do not require teacher intervention and can rather substitute human assessment, was identified as one of the guiding motivations behind the research activities undertaken throughout the PhD.


\subsubsection{Automated Evaluation of UML Class Diagrams}
\label{sec:uml_eval}

Automatic assessment of UML diagrams is a relevant problem for software engineering educators, and research literature has explored multiple approaches to assess it. A valid assessment process cannot simply evaluate whether a diagram follows rules and constraints defined by the UML specification, but should also ensure that the intent behind a model is consistent with the scenario it is supposed to represent, evaluating the students' ability to represent adequately different domains together with following modeling best practices.

One of the earliest studies that systematically defined ways to understand and measure model quality was presented by Lindland et al. \cite{lindland1994understanding}, who identified three quality dimensions: \textit{syntactic quality}, which was defined as the adherence to a given language's modeling rules, \textit{semantic quality}, presented as a diagram's fidelity to the intended domain, and \textit{pragmatic quality}, the understandability of a diagram to human readers and stakeholders. Over the years, these three dimensions became the most commonly used taxonomy to evaluate model-based assignments such as UML diagrams, with several automated tools referring back to them. 

An application of these three dimensions specifically for UML class diagrams was described by Bolloju and Leung \cite{bolloju2006assisting}, who analyzed recurring modeling errors made by students and grouped them according to the three categories. The resulting taxonomy found that most syntax issues regard absence of cardinality values or attribute types, inappropriate element naming, and general incorrect use of the UML notation when representing a diagram; semantic issues related to wrong multiplicity values or attribute types, placing attributes or methods in the wrong classes, and using wrong association types when connecting classes; finally, pragmatic issues covered understandability defects such as redundant elements, under-explained hierarchies in inheritance relationships, readability issues, and inconsistencies in the graphical representation of elements.

Following the taxonomy described above, multiple tools have been developed over the last ten years to automate class diagram evaluation: one example of such tools has been developed by Anas et al. \cite{anas2021newmethod}, who implemented a semi-automatic summative approach for providing feedback to students without relying on teacher intervention. Their solution works by converting a diagram into a labeled graph where classes and attributes are nodes and associations between classes become the connecting edges; a similarity measure computed as the weighted sum of three independent, configurable measures (syntactic, structural, semantic) between said graph and the graph representation of the teacher's diagram is used to generate feedback showing matching elements, differences, and errors. A certain degree of variability is allowed in the valuation by accepting comparisons with multiple teacher diagrams, but the authors highlight the weakness of the approach on complex exercises, where matching performance tends to fall, leading to low precision scores.

Another approach for automated comparison was presented by Nikiforova et al. \cite{nikiforova2015approach}, who defined a numeric metric named \textit{semantic distance} to quantify the difference between two diagrams as a way to understand missing elements and matches without needing to manually inspect them. The approach consists of computing four independent distance vectors for classes and interfaces, attributes, methods, and associations: for each vector, elements of the corresponding type found in the first diagram are paired with the element in the second with the closest semantic similarity (found by comparing the elements' names); the euclidean vectors are then used to compute the semantic distance for each element type, and for the overall diagram, allowing for a way to understand which areas need to be revised to bring a diagram closer to a reference. An overall distance score of 0 means that two diagrams are equivalent, and higher distance scores mean that the diagram under evaluation is farther from the reference; while the approach is effective for understanding improvement areas, it requires human intervention in case of elements with names that are semantically equivalent but lexically different, since they would not be considered a valid match. 

UMLGrade \cite{farias2020whats} is presented as a six-step pipeline where a teacher can define the reference diagram and evaluation criteria, and an automated comparison can then be performed between said reference and student-produced diagrams, resulting in a report that covers semantic and syntax aspects, as well as adherence to design rules and object-oriented development principles, and readability, going beyond simple structural comparison. It can be seen as a way to emulate the evaluation process performed by a teacher, and allows multiple ways to convey the same concept, allowing for more variety in the space of accepted student solutions. 

A practical example of how UMLGrade can be used in practice has been described by Modi et al. \cite{modi2021tool}, who extended it to also cover additional UML notations such as use case, activity, sequence, and state machine diagrams. The tool has been implemented as a desktop architecture where a teacher can upload a zip archive with the student diagrams and the source code of their solution, and different components are used to identify the type of diagram, compare it with the corresponding reference, and generate a feedback report for the students, showing the list of correct elements and missing ones with the overall score, and one for the teacher with the recap of assigned grades. The matching approach, however, does not allow flexibility and expects explicit name equivalence, meaning that semantically equivalent representations are considered wrong; there are no examples of actual tool use with students, and no comparison with human evaluators, meaning that its strengths and potential benefits are only assumed. 

Bian et al. \cite{bian2019automated} identified the problem of variants and alternative representations for UML class diagram assignments, and proposed a metamodel and a three-step grading algorithm to automate diagram evaluation. The first step takes every class in a teacher-defined reference diagram and looks for the closest class in the student's diagram, considering syntax similarity, semantic similarity using techniques such as the WordNet database \cite{miller1995wordnet} or the identification of semantic relationships among names, and structural similarity (similarity of attributes and methods) to identify the best match; following steps match attributes and methods in matched classes using the same approach, then match associations between matched classes. A preliminary study performed with a tool implementing this approach yielded promising results (a grade distribution within 14\% of the human-assigned grades), although no examples of further use exist in the literature.

Lastly, Wodel-EDU \cite{gomezabajo2017dsl} implements a different approach: rather than automatically evaluating student-produced diagrams compared to a reference, it produces variants, or \textit{mutants}, of said reference and asks students to identify the correct model. This is achieved through the domain specific language (DSL) Wodel \cite{gomezabajo2016wodel}, which can be applied to any type of modeling language according to a changeable metamodel and allows for multiple operations such as creating, editing, or deleting objects with multiple selection criteria and support for computing static and dynamic element coverage. 

Most of the existing solutions implement name-matching techniques to identify correct modeling choices by computing string similarity between named elements, a measure that can help identify matching elements while also allowing sufficient lexical variation (e.g., typos) between equivalent concept labels. An extension of this concept, named the normalized edit distance, or Levenshtein distance, is used to quantify the minimum number of single-character changes (additions, deletions, substitutions) that are necessary to convert one string into another, normalized to fall within the [0, 1] interval so that distance scores below a given threshold can be considered matches independently of the length of the two compared strings \cite{navarro2001guided}.

\subsubsection{Automated Generation of UML Class Diagrams}

In addition to the tools and approaches described in the previous subsection for automated class diagram assessments, the related research literature also presented approaches for automatically generating such types of diagrams from natural language requirements via multiple techniques, showing relatively limited success. 

Visual Narrator \cite{robeer2016automated} is one example of these tools, extracting conceptual models from user stories (requirements expressed in the format \textit{As an <actor> I want to <goal> so that I can <gain>}) through a pipeline of eleven natural language processing heuristics such as part-of-speech tagging to identify nouns, verbs, and adjectives, chunking of phrases, and resolution of common references. The tool was tested on two case studies, achieving precision and recall scores between 80\% and 92\%, indicating the potential effectiveness of the tool.

A similar approach was presented by Yang and Sahraoui \cite{yang2022towards}, who proposed a grammar-based pipeline to extract class diagram from specifications expressed in natural language with sentence classification being used as the basis for the generation of diagram fragments; however, the limited precision and recall scores achieved on the dataset used to evaluate the solution led the authors to conclude that rule-based methods do not have enough semantic flexibility to allow for complete specification of the requirements.

A machine learning reliant approach was instead introduced by Saini et al. \cite{saini2022machine}, whose incremental and interactive tool for domain modeling allowed humans to partake in the requirements extraction process, correcting the produced artifacts and demonstrating that a fully automated approach is less effective compared to an iterative process that supports automation with human oversight.

Other tools were developed with a focus on educational use, allowing for the generation of UML diagram exercises. One such tool is AutoER \cite{foss2022automatic}, that supports automatic generation of UML diagrams and subsequent variants through the randomization of element names and relationship parameters. These variants can then be used to compose exercises where students can mark the diagrams and receive feedback by interacting with the exercise's accompanying textual description. The Wodel-EDU framework was also integrated in the online educational platform Moodle \cite{gomez2023automated}, integrating seven types of domain-independent exercises automatically generated with the mutants, allowing for automated grading and feedback .

While educational tools to support UML education are common, the majority of them support mostly class and sequence diagrams, while use case, activity, and state machine diagrams remain relatively less covered \cite{huber2022tool}; another gap in these types of tools is that the tools that automate the exercise generation process rarely tend to include formative feedback that depends on the students' modeling choices, highlighting the need for platforms that include exercise creation for teachers, student diagram editing in free-form, and automated evaluation as functionalities served together.

\subsubsection{Business Process Modeling and BPMN}
\label{sec:back_bpmn}

The Business Process Model and Notation (BPMN) is a graphical standard used to standardize the representation of business processes by representing its different elements in order to capture and define activities, rules, actors, and responsibilities involved in the process. The elements that compose a BPMN diagram belong to five distinct categories:

\textbf{Tasks} are also known as Activities and represent the different actions that are performed during a business process and can be distinguished by their type, with \emph{User Tasks} (activities performed by a human actor by using a software application), \emph{Service Tasks} (automated activities performed by software applications), and \emph{Manual Tasks} (activities performed by a human actor without any software intervention) being the three most commonly used ones.

Another important feature offered by Tasks is the fact that they can be used to represent \emph{Sub-Processes}, that is, groupings of multiple elements in sequence that can enable process decomposition effectively and improve the organization and readability of the process model. Tasks are represented as rectangles with curved angles, and can contain a certain icon that distinguishes their type (e.g. a human icon, a hand, gears, a letter, or a plus shape for sub-processes).

\textbf{Flows} are arrows that are used to connect the different elements in a BPMN diagram; the most common and important type is the \emph{Sequence Flow}, which is used to display how elements are sequentially connected (e.g., the element a sequence flow is going out of comes before the element touched by a flow's arrow), while another relevant type is the \emph{Message Flow}, which displays how elements in separate processes exchange information through messages.

\textbf{Events} represent something that can happen through a process without being directly performed by a human actor and impact how it continues, and can be used to begin a process (\emph{Start Events}), impact it during its execution (\emph{Intermediate Events}), or conclude it (\emph{End Events}). They are drawn as a circle and represent various scenarios such as receiving/sending a message to another process, a certain time condition (e.g., a date, the passage of a specific interval of time) becoming true, or internal server errors.

\textbf{Gateways} control the sequence flows inside processes by allowing branching or joining operations; gateways are expected to be represented as pairs, with one having multiple outgoing flows acting as the splitting gateway, followed later in the process by another one of the same type that merges back the branched incoming flows. They can be \emph{Exclusive} (if their outgoing flows are associated with conditions, and only one condition can be true at one time, signifying which flow can proceed), \emph{Parallel} (if all outgoing flows are allowed to proceed at the same time, with no condition associated), \emph{Inclusive} (if their outgoing flows are associated with conditions, and one or more conditions can be true at the same time), or \emph{Event-Based} (if their outgoing flows are followed by events, in which case only one event can have its associated condition become true, specifying which flow can proceed).

\textbf{Pools} are rectangle-shaped containers that represent the different business entities that participate in a process; they can be divided into \emph{Lanes} that specify the various human actors that belong to the same entity and are involved in the process. Pools can also be represented as \emph{Collapsed} if they represent an entity that is involved with the process but is external and thus no information is known about it; communication between different pools and elements inside them is allowed exclusively with message flows.

Despite its industrial relevance, BPMN is recognized as a challenging notation for novice learners, due to recurring obstacles such as the absence of standardized pedagogical frameworks, the varying student preparation levels affecting the learning process and, most importantly, the cognitive difficulty that comes from the need to map real-world process descriptions onto a formal graphical representation \cite{bandara2010bpm}. Novice modelers also tend to fall into common error patterns related to correct control flow representation, handling of events, and communication between separate pools, requiring continuous practice to reach adequate modeling proficiency \cite{maslov2022towards}. These common error patterns and obstacles tend to make BPMN modeling frustrating and anxiety-inducing, and such negative effects have been observed to have a significant impact on the quality of the produced models, as well as on the satisfaction of the learning process \cite{Nobre2025ModelQA}.

Another obstacle related to student BPMN appreciation comes from the difficulty in assessing whether a given diagram is correct, a challenge that also involves lecturers and practitioners. The standard reference for BPMN quality was introduced by Dumas et al. \cite{dumas2018fundamentals} and identifies three quality dimensions: syntactic quality, which measures how much a given diagram respects the BPMN standard's formal grammar and structural rules, semantic quality, the evaluation of how well a diagram represents the process it is intended to capture, and pragmatic quality, which assesses how readable and understandable a diagram is for the intended target audience (e.g., stakeholders in industrial projects). Another point of reference consists of the Seven Process Modeling Guidelines \cite{mendling2010seven}, a set of empirically grounded rules to apply in the creation or evaluation of a process model to ensure quality measures are respected such as using the minimum set of elements possible to convey a given requirement, or following naming conventions when labeling elements.

Given the large set of criteria tied to assessing BPMN diagrams' correctness, automated verification of such models has emerged as an active research area, driven by the necessity of providing accurate and consistent feedback. Lopes and Guerrero \cite{lopes2023assessing} identified a set of formal verification and testing techniques that can be used to evaluate BPMN diagrams such as model checking, black box or regression testing, and methods to verify constraint satisfaction, while also proposing a framework to structure research steps on process model evaluation. However, the approaches found by the two authors' review are mostly focused on industrial process validation and, while they are effective in identifying structural and behavioral errors, they do not provide feedback at the element level, or consider the domain to be represented when assessing models, making them less suitable for educational environments where students need to identify their mistakes and know how to correct them.

%G9 is not really necessary, the thesis seems to be fine without this paragraph
%BPMN education involves practical assignments where students are tasked with representing a given process description in graphical form; such activity relies on dedicated software environments that let students draw and edit diagrams with ease, facilitating the production process without relying on pen and paper drawing. One of the most commonly used web-based environments in university courses is the Signavio Academic platform, initiated as the Signavio-Oryx Academic Initiative \cite{kunze2010signavio}: what makes the platform particularly effective is its support for multiple modeling notations commonly used in industrial settings, including BPMN, UML, and Petri nets, as well as allowing multiple users to perform collaborative editing of their diagrams and the use of dedicated submission workflows. Alternative platforms do exist, such as the open-source \textit{bpmn-js} library\footnote{\url{https://bpmn.io/toolkit/bpmn-js/}} 

\section{Gamification}

Gamification is a concept that has been formally defined in 2011 by Deterding et al. \cite{deterding_game_2011} as \emph{the use of game design elements in non-recreational contexts}; its intended goal is to enhance activities that are considered unappealing, unpleasant, or uninteresting by increasing user motivation and interest, with improvements in performance being another possible benefit.

An important condition to be observed when defining gamification is to distinguish between a \textit{gamified system} and a \textit{serious game}. A serious game is defined as a completely interactive experience (a video game, a board game, or a simulation environment) whose primary objective is not entertainment, which is a secondary consideration, but rather learning, training, or persuasion are seen as the main focus \cite{michael2006serious}. In contrast, gamified systems do not create a game but extract game design elements and combine them with an existing activity, application, or environment with a non-recreational purpose. The distinction implies that a serious game requires learners to engage with the game mechanics to learn while a gamified system augments a task that can be performed normally with motivational mechanics while leaving the original activity unchanged.

The formal definition of gamification is purposely agnostic in relation of which elements should be applied, given that the different elements may have varying effects and may thus be better suited for specific contexts. The categorization by Werbach and Hunter \cite{werbach2012win} identified a three-tiered hierarchy, distinguishing game mechanics between: I) components, concrete, tangible elements such as points, badges, leaderboards; II) mechanics, rules that guide player behavior such as challenges, competition, or cooperation; III) dynamics, system-level phenomena that support motivational aspects such as emotions, relationships, progression, and emerge better over continuous use. Thanks to the hierarchy, it is possible to identify which game mechanics depend on a surface level with little reasoning (e.g., adding a point-based leaderboard), and which are implemented in educational aspects as a consequence of a motivational design.

It must be observed, however, that not all game elements can be equally effective in practice: a meta-analysis of empirical studies adopting gamification \cite{hamari2014does} showed that, while gamification leads to positive outcomes in the majority of cases, the scale of positive effects strongly depends on both the implementation context and the characteristics of the intended users; for example, leaderboards and other competitive elements may alienate certain users and make them feel less motivated to continue interacting with gamified tools.

To understand the psychological foundations behind gamification, it is necessary to distinguish between intrinsic and extrinsic motivation. Intrinsic motivation is agreed upon as the act of performing activities for inherent satisfaction, where individuals derive enjoyment from the processes they are undertaking and feel motivated to take on new challenges for the sake of it; extrinsic motivation is instead considered as the behavior driven by factors that are outside of one's control such as material rewards, social recognition, or even negative effects and punishments that increase pressure to complete certain activities.  

Motivational aspects derived from this distinction are usually explained according to Self Determination Theory (SDT) \cite{deci2000intrinsic, deci1985intrinsic}: intrinsic motivation relies on satisfying three basic psychological needs such as competence (feeling capable and observing the effects of one's actions), autonomy (experiencing choices and being allowed to have multiple directions and options to perform tasks), and relatedness (being connected to other people). Gamification approaches should be designed to cover all three needs with carefully selected mechanics: the use of progression mechanics such as experience points and progress bars support competence and make one's growth easy to observe; avatars, storytelling, and customization options allow one to express themselves autonomously; cooperative and competitive mechanisms such as leaderboards and group-based tasks implement relatedness to others. Intrinsic motivators are necessary, and should be the guiding force behind the design of gamified tools and approaches: if game elements rely on extrinsic motivators such as points or rewards users may feel less inclined to interact with the game environment over time, leading to progressive loss of motivation.

Among the several frameworks that have emerged over the years to categorize game mechanics and guide the creation of gamified solution, Octalysis \cite{chou2015actionable} is particularly remarkable due to how it identifies the three needs with intrinsic and extrinsic motivators and sorts them according to eight distinct \textit{Core Drives}, dimensions under which game mechanics can be categorized depending on whether they bring intrinsic or extrinsic motivators, as well as whether their motivational effect is derived from positive or negative impact; for example, rewards for correct actions, social recognition, feedback after one's choices are positive motivators, while penalties in case of errors, resource scarcity, or limited content are negative motivators. The eight Core Drives are summarized in Table \ref{tab:octalysis_core_drives}, with a brief description of each, together with examples of corresponding game mechanics and an indication of whether the drive is a positive or negative motivator, as well as an intrinsic or extrinsic one.

\begin{table*}[ht]
\centering
\footnotesize
\caption{Core Drives defined in the Octalysis framework, with motivational classification and example mechanics.}
\label{tab:octalysis_core_drives}
\resizebox{\linewidth}{!}{%
\begin{tabular}{@{}p{3.2cm}p{6cm}p{2cm}p{2cm}p{4.5cm}@{}}
\toprule
\textbf{Core Drive} & \textbf{Description} & \textbf{Intrinsic /} \newline \textbf{Extrinsic} & \textbf{Positive /} \newline \textbf{Negative} & \textbf{Example Mechanics} \\
\midrule

\emph{Epic Meaning} \newline \emph{and Calling} &
The feeling that one's actions contribute to something greater; feeling part of a broader, shared mission; feeling to be contributing to a common goal. &
Intrinsic & Positive &
Lore and narrative, community-driven quests, collaborative world-building \\
\addlinespace

\emph{Accomplishment} &
Being guided by the desire to make progress, develop skills, and overcome meaningful challenges; feeling the significance of rewards only after earning them through genuine effort. &
Intrinsic & Positive &
Experience points, achievement badges, skill trees, progress bars \\
\addlinespace

\emph{Empowerment} \newline \emph{of Creativity} \newline \emph{and Feedback} &
Deriving motivation by being able to express oneself, to experiment with different strategies and combinations of actions, and to receive meaningful feedback that validates one's actions. &
Intrinsic & Positive &
Open-ended challenges, immediate feedback, customization options, sandbox modes \\
\addlinespace

\emph{Ownership} \newline \emph{and Possession} &
Feeling ownership over virtual or tangible goods, with the consequent drive to protect and improve what is owned, accumulate new items, and invest more time in the system. &
Extrinsic & Positive &
Virtual currencies, avatar customization, inventory systems, collectibles \\
\addlinespace

\emph{Social Influence} \newline \emph{and Relatedness} &
Adhering to social dynamics such as acceptance, companionship, competition, and envy; feeling driven to act after observing peers achieve goals or acquire rare rewards. Certain components such as companionship are intrinsic, while competition and comparison mechanics tend to be more extrinsic in nature. &
Intrinsic \& Extrinsic & Positive &
Leaderboards, cooperative tasks, guilds and teams, social recognition \\
\addlinespace

\emph{Scarcity} \newline \emph{and Impatience} &
Being drawn to elements that are scarce or difficult to obtain; drawing interest in something that cannot be obtained immediately and putting sustained effort and commitment in a system as a consequence. &
Extrinsic & Negative &
Limited-time rewards, locked content, appointment mechanics, exclusive items \\
\addlinespace

\emph{Unpredictability} \newline \emph{and Curiosity} &
Feeling engaged in a system thanks to the desire to discover what comes next; being stimulated by varied and unpredictable experiences that reduce monotony and sustain long-term interest. &
Extrinsic & Positive &
Random rewards, mystery boxes, surprise events, Easter eggs \\
\addlinespace

\emph{Loss} \newline \emph{and Avoidance} &
Acting carefully and promptly in order to avoid negative consequences such as losing something already acquired, or of missing a temporary opportunity. &
Extrinsic & Negative &
Penalties, health/life systems, expiring rewards, failure states \\

\bottomrule
\end{tabular}%
}
\end{table*}





\subsection{Gamification in Educational Contexts}

The use of gamification for educational support has been extensively studied ever since its formal definition, leading to the creation of a growing body of empirical evidence on both its potential effectiveness and its limitations. A systematic mapping over 69 publications \cite{dicheva2015gamification} identified the most commonly used elements in the form of points, badges, and leaderboards, found that the educational domains where gamification is applied the most are computer science and programming education, and, most importantly, identified two key limitations: most of the studies analyzed lacked rigorous empirical evaluation and relied more on anecdotal evidence, with few controlled experiments to support claims; another relevant issue was the absence of theoretical grounding, since very few studies justified design choices with motivational theory or validated gamification frameworks.

The literature review by Hamari et al. \cite{hamari2014does} addressed the effectiveness of gamification in the educational domain, reviewing specifically empirical studies that used it as a strategy: most of the educational studies reported positive outcomes in terms of student engagement and participation; however, a key limitation had emerged in relation to the dependency on the context in which gamification is used, which in turn may lead to negative or neutral results if the design of a gamified tool is not built adequately (e.g., an excessive focus on extrinsic motivators or rewards for the sake of rewarding). 

One of the earliest examples of gamification being used in engineering courses showed that the adoption of experience points, levels, badges, and challenges led to significant increases in lecture attendance, participation on the course's online platform, and students reported feeling more motivated, finding the course more engaging and interesting compared to traditional lectures \cite{barata2013engaging}. 

However, other early studies showed more mixed results \cite{dominguez2013gamifying}: while the use of a gamified plugin for an e-learning platform led to higher grades for practical assignment and end-of-course evaluations, students who used the plugin performed worse on written tasks and participated less actively during lectures, suggesting that game mechanics may shift the focus from learning to simply accumulating points and rewards. A longitudinal study by Hanus and Fox \cite{hanus2015assessing} yielded similar results: the comparison of a course with leaderboards and badges as motivating elements with a non-gamified equivalent course over 16 weeks found that students in the gamified course eventually showed less intrinsic motivation, satisfaction, and empowerment, participating less in activities, and eventually achieving lower grades in the final exam. The two studies support the importance of designing gamified educational tools with care, and suggest that adding competitive mechanisms may not be beneficial in all scenarios.

A systematic mapping on the negative effects of gamification \cite{toda2017dark} found that there are four main categories to consider: the reduction of performance, the gradual loss of motivation, undesired behaviors, and student indifference, with causes behind these issues being inadequate planning of the gamified experience, excessive competition, penalization, and long-term use that makes students feel less interested in interacting with the gamified solutions. The authors then followed up on their review with a taxonomy of game elements across five categories (Performance, Social, Personal, Ecological, Fictional) to support practitioners in the selection of game elements aligned with their specific goals instead of choosing commonly used mechanics \cite{toda2019taxonomy}.

Automated and immediate feedback is considered one of the most critical choices for educational gamification according to formative assessment literature. Good feedback practice, according to Nicol and Macfarlane-Dick \cite{nicol2006formative}, should be structured around seven principles that include timeliness, specificity, and the support for self-regulation: feedback should be received promptly, target identifiable errors, and provide actionable guidance to correct the errors' causes, allowing students to iteratively improve their work before a summative evaluation by lecturers. The seven principles also distinguish between formative feedback, which is provided during the learning process and guides improvement, and summative assessment that evaluates students after learning is completed. Gamified systems should follow these principles and aim to tie automated feedback directly to the implemented game mechanics, to provide a more complete and effective learning process.

A more recent meta-analysis focused on gamification's educational effectiveness \cite{sailer2020gamification} found that it can have small or medium positive effects on cognitive learning outcomes, motivation, and behavioral engagement; it is worth mentioning that cognitive improvements were the most commonly observed, while motivational and behavioral effects showed higher variability, hinting that the design quality of gamified systems impacts the width and duration of the benefits. The findings build upon the review by Hamari et al. \cite{hamari2014does} and reinforce the conclusion that gamification can positively impact learning processes as long as mechanics are carefully matched to the educational activity and to the needs of the student population.

A challenge to consider in gamification research is the so-called novelty effect, where initial improvements in student engagement may be caused by the novelty of the environment itself rather than from the intrinsic quality of the game mechanics, with the risk that such gains may stabilize or even fall when students become familiar with the system. Longitudinal studies have experienced such an effect: across a fourteen-week semester with gamified enhancements, the novelty effect took place after approximately four weeks, leading to a loss of motivation for most of the remaining part of the semester, until students adapted to the system with a familiarization effect and showed renewed motivation to interact with the gamified aspects \cite{rodrigues2022gamification}. The authors argued that more complex gamification mechanics built on cooperation, competition, and narrative events were considered more novel and interesting than simpler mechanics such as points and badges, indicating that the longitudinal use of gamified tools should be designed with more advanced mechanics as the key focus to sustain long-term engagement.



\subsection{Gamification in Software Modeling}
\label{sec:back_gamif_mde}

While gamification has been adopted in multiple ways in educational contexts, examples of its application focused specifically on software modeling activities are somewhat fragmented. The two modeling notations covered by this PhD thesis (BPMN, UML class diagrams) have attracted a growing body of attention from the research community, reflecting the educational importance of these topics together with the need to make their learning more engaging.

Examples of gamification or game-like approaches being used specifically in relation to BPMN modeling are not common in literature: the only example of gamification being used for BPMN is BPMS-Game \cite{mancebo_bpms-game_2017}, a tool that is used in an industrial context to encourage modeling according to good sustainability practices. The rules can be defined by administrators, and the game mechanics employed by the game include badges, achievements, and leaderboard; the tool is not intended for educational use, however, and is more focused toward practical use, making it unsuitable for classroom use.

Other solutions do exist, but are serious games rather than gamified tools: one such game was described by Marín \cite{marin_lessons_2022}, and gives students a set of challenges with three increasing difficulty levels (basic, medium, advanced); students are given an incomplete process and have to choose the elements needed to complete it out of a set of given ones, and theoretical questions are also included with the exercises to test the students' knowledge of good modeling practices. A score is computed depending on how many elements were used, how many correct answers were given, how much time the student needed, and how correct the resulting diagram was; this score is then used for a leaderboard where students can rank themselves with their peers, and badges can be obtained in case of good performance (e.g., correctly answering all questions in a level, achieving 100\% correctness in an exercise).

Another serious game, that puts more emphasis on competition between groups of students, is BPMN-Wheel \cite{kutun_bpmn_2019}, a board game where two teams of players have to answer theory questions on BPMN modeling practices to gain either in-game currency or elements that can be used to compose a process that has to be represented; the team whose resulting process is the most similar to the target one wins the game.

A broader set of tools has been proposed for UML class diagram education, with both serious games and gamified approaches being developed and discussed as potential solutions to enhance the learning process and make modeling activities more attractive for students. A mapping review, which is discussed more in detail in Section \ref{sec:back4}, has been performed to identify the set of relevant studies that discuss game-like approaches that can be applied to conceptual modeling education: following the retrieval of all sources, each study was analyzed to identify the game mechanics it implemented, the adherence to the Octalysis framework in terms of which Core Drives were being followed in the design and, finally, whether the papers described a gamified tool or a serious game.

Jurgelaitis et al. \cite{jurgelaitis2019implementing} described the use of gamification elements such as points, achievements, leaderboards, and badges to teach UML in a level-based course structure using an online platform where students were asked to answer multiple choice questions on different diagram types; unlockable example diagrams were also implemented with the virtual currency used to unlock them being tied to exercise completion. Results showed that the mechanics improved both the average student scores and their self-declared intrinsic motivation.

A similar approach was implemented with LOOP \cite{krugel2017computational}, an online course focused on object-oriented programming with a block-based editor to draw class diagrams, offering feedback for students as well as visual changes in the application when changing the diagram. Similarly, Denden et al. \cite{denden2021effects} implemented gamification elements such as bonus points that can be obtained for completing certain assignments, badges that reward completion of quizzes on specific course topics, and leaderboards to rank students according to their experience points, with UML class diagrams being among the topics covered by the course. 


ModelGame \cite{junior2021modelgame} is described as a reference framework that teachers can use to evaluate UML diagrams according to a set of parameters, represented by ten different quality metrics such as scope, syntactic, semantic, detection of inconsistencies, and resolution; it makes use of missions as a way to motivate students, with progression, points, and feedback being tied to adhering to the quality metrics.

A serious game described by Livovský and Porubän \cite{livovsky2014learning} was presented as a 2-dimensional adventure where levels, challenges, and enemies were tied to object-oriented topics through an Object Access Tool used to connect game elements to class diagrams. Modifying the diagram would create new elements in the game, or change the properties of existing ones, while interacting with certain game objects would highlight the corresponding class and activate its methods.

Another serious game is Classutopia \cite{marin2018learning, larenas2018classutopia}, which consists of a role-playing game that has students act as a robot battling against an evil wizard in a ruined future. Students can interact with class diagrams to customize the robotic avatar's features and abilities, and solving defective diagrams is the main gameplay loop: making changes that improve a diagram's correctness damages the enemy, while errors sabotage the player, and diagram coloring is used to highlight both correct changes and mistakes.

GaMoVR \cite{yigitbas2022gamification, yigitbas2023design} is an innovative gamified solution that has students solve exercises in a virtual reality environment where it is possible to rearrange and connect elements directly. The tool offers different modes of play: one where requirements are given and a solution must be created, with a Hangman game tied to the number of errors made by the student, and another mode where students have to correct a diagram with errors by shooting at the incorrect elements. Gamification elements also include experience points and levels that unlock additional tools and bonuses.

innov8 \cite{boughzala2017feedback} is a serious game that has students play the role of business consultants for a fictional company, with the goal of creating the class diagram that represents a certain set of requirements, gathering the necessary information through playing the game and interacting with non-playable characters. Students can then define a collective solution and compare it to the reference drawn by the teacher, receiving feedback after a comparison. 

Dæhli et al presented LearnER \cite{daehli2021exploring}, a gamified tool that supports Entity-Relationship models in addition to UML, and works by having students select elements out of a set of alternatives to compose the model that best represents the domain described in a textual assignment. Game mechanics include points awarded for completing exercises, hints that can either be obtained by choosing correct elements or bought by reducing the amount of reward points, and leaderboards for the number of points obtained, both overall and for each exercise.

A serious game-based approach was described by Dupuy-Chessa et al. \cite{dupuychessa2015lego}, who applied the concepts of the game Pictionary to diagram modeling: with the game, two teams of students face each other on a board game where cards can be selected; these cards specify the requirements that students have to draw in a limited time window, with teacher feedback being an available option.

Lastly, Sanjaya et al. \cite{sanjaya2020development} theorized a gamified mobile application for gamifying class diagrams, among other software engineering topics, with interactive quizzes focused on system design and requirements definition, using badges, achievements, and rankings as motivational elements.

Table \ref{tab:mapping} summarizes the literature mapping by listing all tools and approaches found and, for each one, whether it is a gamified tool or a serious game, which game mechanics are implemented, and which Core Drives are followed.

{\begin{table*}[ht]
\centering
\footnotesize
\begin{tabular}{@{}p{3cm}p{2.5cm}p{4cm}p{4cm}@{}}
\toprule
\textbf{Tool Name} & \textbf{Strategy} & \textbf{Game Mechanics} & \textbf{Octalysis Core Drives} \\
\midrule
GaMoVR \cite{yigitbas2022gamification, yigitbas2023design} & Gamification & Levels, Points, Unlockable Elements, Feedback & Accomplishment, Loss and Avoidance \\ 
Moodle course (Jurgelaitis et al.) \cite{jurgelaitis2019implementing} & Gamification & Points, virtual currency, leaderboard, badges & Accomplishment, Ownership \\ 
LOOP\cite{krugel2017computational} & Serious Game & Feedback & Accomplishment \\ 
Moodle course (Denden et al) \cite{denden2021effects} & Gamification & Points, badges, leaderboards & Accomplishment
Gamified mobile app \cite{sanjaya2020development} & Gamification & Challenges, badges, leaderboards & Accomplishment \\ 
LearnER \cite{daehli2021exploring} & Serious Game & Progress, points, hints, feedback, leaderboards & Accomplishment \\ 
ModelGame \cite{junior2021modelgame} & Gamification & Points, progress, feedback & Accomplishment \\ 
innov8 \cite{boughzala2017feedback} & Serious Game & Cooperation & Social Influence \\ 
Pictionary-like \cite{dupuychessa2015lego} & Serious Game & Cooperation, feedback, rewards & Social Influence, Scarcity, Accomplishment \\ 
Classutopia \cite{marin2018learning, larenas2018classutopia} & Serious Game & Feedback, Challenges, Unlockable Rewards & Epic Meaning, Accomplishment, Empowerment \\ 
2D action-adventure game \cite{livovsky2014learning} & Serious Game & Levels, Feedback & Epic Meaning, Accomplishment, Empowerment, Unpredictability \\
\bottomrule
\end{tabular}
\caption{Summary of existing game-like solutions for UML class diagram education}
\label{tab:mapping}
\end{table*}

The mapping process was conducted in early 2023, following the first empirical study described in Chapter \ref{ch3}, with the goal of identifying existing tools for UML class diagram education. The process was then repeated in late 2024 to retrieve additional studies that had emerged after the conduction of the UMLegend controlled study and observe evolutions in the landscape of gamified modeling tools. Three new tools emerged as a result of the second analysis; the summary of the mapping analysis executed on them is reported in Table \ref{tab:mapping2}.

Among the three tools, Papygame \cite{bucchiarone2020papyrus, bucchiarone2023gamifying} was described as a gamified plugin for the professional modeling environment Papyrus, which supports multiple UML notations. The main feature of the plugin is the structure of UML assignments as a sequential challenge with increasing difficulty level, exploiting the Hangman game to represent the modeling challenge: students can select diagram elements to compose a solution, with each incorrect option adding a new part of the hanged drawing, with immediate feedback displaying the errors. Successful exercise completion unlocks new challenges and awards experience and virtual currency. Experiments using the tool showed that it can have positive effects on productivity and learning outcomes.

The same authors also integrated a 2-dimensional virtual environment with the Papyrus web editor to implement PolyGloT-UML \cite{bucchiarone2024polyglot}, an educational platform that supports learning paths that students can customize according to their progress, teacher-defined activities such as quizzes or exercises, and assignments that award badges; most importantly, the tool is the only example of gamified modeling tools that implements social dynamics, allowing students to interact with each other in the virtual hub. 

Lastly, ModelDefenders \cite{cammaerts2023modeldefenders, cammaerts2024towards} adapts the concept of mutation testing to Model-Driven Engineering: students receive a class diagram and can face one another as either defenders who define action sequences based on the model and their expected output, or as attackers that apply changes to the diagram to introduce bugs; mutants are then considered killed if there is at least one test case that has a different output on the mutated diagram compared to the original one.

{\begin{table*}[ht]
\centering
\footnotesize
\begin{tabular}{@{}p{3cm}p{2.5cm}p{4cm}p{4cm}@{}}
\toprule
\textbf{Tool Name} & \textbf{Strategy} & \textbf{Game Mechanics} & \textbf{Octalysis Core Drives} \\
\midrule
Papygame \cite{bucchiarone2020papyrus, bucchiarone2023gamifying} & Gamification & Points, Feedback, Progress, Levels & Accomplishment, Loss and Avoidance, Empowerment \\ 
PolyGloT-UML \cite{bucchiarone2024polyglot} & Gamification & Feedback, challenges, socialization, badges & Accomplishment, Ownership, Social Influence \\ 
ModelDefenders \cite{cammaerts2023modeldefenders, cammaerts2024towards} & Gamification & Points, competition, feedback, challenges & Accomplishment, Empowerment, Social Influence \\
\bottomrule
\end{tabular}
\caption{Summary of game-like solutions for UML class diagram education found with the second literature mapping analysis}
\label{tab:mapping2}
\end{table*}

Taken together, the landscape of gamified tools for both BPMN and UML class diagram education reveals a common pattern: existing solutions focus mostly on single-session use and tend to rely on the selection of answers out of predefined element sets instead of allowing for free-form diagram construction. The consequent research gap is the need of tools that combine a free modeling environment with automated, gamification-based feedback.

\section{Artificial Intelligence and Large Language Models}

Large Language Models (LLMs) are neural sequence-to-sequence models built upon the Transformer architecture \cite{vaswani2017attention}, and are considered a significant advancement on neural architectures thanks to self-attention mechanisms and positional encoding, allowing for easy training on massive textual datasets and superior performance on text processing; LLMs are trained to be text predictors that divide text into tokens and, depending on their training data, produce compatible output depending on the input they receive, which is conventionally named a \textit{prompt}; these architecture can be trained to be used for general-purpose tasks or fine-tuned so that they are better aligned to a specific goal.

The most influential family of LLMs, as well as the one synonymous with artificial intelligence for everyday use, is the GPT series developed by OpenAI: the first model to be made available for public use through a chat interface (the so-called ChatGPT) was GPT-3 at the end of 2022, and demonstrated that scaling the size of a model to 175 billion parameters for training shows relevant capabilities related to instruction following, learning from the context of a prompt and the examples it contains, and structured text generation \cite{brown2020language}. The GPT-4 model released in 2023 \cite{openai2023gpt4} extended these capabilities even further, with improvements such as more precise reasoning, longer context windows (the length, measured in tokens, of messages that can be parsed by the model at a given time), and multimodal support, allowing for input, as well as the production of answers, in non-textual form (e.g., images).

Following the release of ChatGPT, the LLM landscape expanded in 2023-2024 with multiple competing model families emerging such as Gemini \cite{anil2023gemini}, released by Google DeepMind as a set of multimodal agents whose strengths lay in structured text generation and contextual understanding, Meta's LLaMA \cite{touvron2023llama}, a family of publicly available models that allowed for the birth of fine-tuned derivatives for research use, Qwen \cite{bai2023qwen} by the Alibaba Group, focused on strong code generation capabilities, and DeepSeek \cite{deepseek2024}, a family of open source models optimized for code and reasoning that showed competitive performance in comparison to the existing proprietary models.

Given their novelty at the time the PhD was being conducted, multiple surveys explored the application of LLMs to software engineering tasks. A 2023 survey on the emerging research works of LLMs for software engineering \cite{fan2023large} found applications for coding, documentation, software testing, design and requirements activities, and highlighted that LLM capabilities such as creativity and ability to understand and follow instructions could allow for automation opportunities in multiple phases of the software lifecycle, while also pointing out the risks tied to hallucinations (the scenario where an LLM's textual output is nonsensical or inconsistent with the input received, while giving the impression of being natural and fluent despite being incorrect \cite{ji2023survey}) that require careful human intervention and validation.

The risks associated with LLM hallucinations can be characterized in two types, according to the survey by Huang et al. \cite{huang2025hallucination}: factuality hallucinations happen when the model's output conflicts with verifiable real-life knowledge, while faithfulness hallucinations are scenarios in which the answer is inconsistent with the given input prompt by either not following the provided instructions, not respecting the contextual information, or not adhering to the internal logical structure; faithfulness hallucinations are particularly critical in structured generation tasks, since a model may produce syntactically valid answers whose content was never stated or implied by the input fed to the model. Both types of hallucinations require detection strategies such as model-based verification of the output or pipelines that parse the model according to external knowledge, as well as mitigation strategies that include retrieval-augmented generation, iterative self-refinement loops within the model, or multi-agent architectures that continuously interact to refine the final output. 

A literature review by Hou et al. \cite{hou2023large} analyzed 395 studies and found that, at the beginning of 2024, most empirical work focused on coding tasks, while requirements engineering activities, as well as model-based tasks, were comparatively less common. Multiple investigations were also performed in the early years of LLM diffusion: for example, Ma et al. \cite{ma2023scope} explored the scope of ChatGPT's capabilities for code generation, code review, and program repair, finding strong syntax-level capabilities, as well as significant limitations for more complex tasks and security vulnerabilities in a sizable amount of the generated programs. Other early studies reinforced the importance of evaluating LLM output with attention to semantic correctness according to domain correctness \cite{ozkaya2023application}, and also pointed out that LLMs exhibit the ability to identify core domain entities from textual description, while also pointing out the importance of human oversight and verification \cite{belzner2023large}. 

The effectiveness of LLMs also relies on the quality of the prompts used to guide their generation tasks; for this reason, advanced techniques known as \textit{prompt engineering} emerged over time to guide research methodologies beyond unstructured questions. One such technique is called Chain-of-Thought (CoT) prompting \cite{wei2022chain}, and consists of asking an agent to reason step by step before it writes its final answer, eliciting more consistent and structured outputs; this technique is particularly effective in tasks where sequential reasoning is necessary and single-shot prompting may not be sufficient to capture implicit requirements. 

Other prompt engineering techniques differ from CoT prompting in terms of how many details are provided to the model at inference time \cite{sahoo2024prompting}. For example, \textit{zero-shot prompting} is when a model receives a description of its intended task with no examples of how its output is supposed to be structured, or how the data it receives as input is formatted; as a consequence, the model needs to interpret the task and produce a response according to the knowledge it received in its training process. This technique can be extended with \textit{few-shot prompting}, a core capability of GPT-3 \cite{brown2020language}, where a small number of examples are introduced in the prompt, improving the model's performance on structured tasks without needing task-specific finetuning.

Another technique is \textit{role-based prompting}, also known as system prompting or persona assignment: with this type of prompts, users instruct the model to act as if they have a specific identity (e.g., researchers, teachers, domain experts) before it begins working on the task; this technique has been proven to be particularly effective, since it leads the model to focus on terms, knowledge, and constraints that depend on the target domain, making it suitable for structured tasks where the output must be expressed in a specific format \cite{sahoo2024prompting}.


\subsection{Large Language Models for Software Modeling}

The use of LLMs to support model-based tasks such as diagram generation emerged as a topic in 2023, with earlier reviews \cite{fan2023large, hou2023large} pointing out that diagram-based tasks were far less common in empirical studies compared to code-focused activities. 

One of the earliest empirical studies on the topic was conducted by Arulmohan et al. \cite{arulmohan2023extracting}, who compared ChatGPT against an existing, state of the art NLP tool on extracting models from domain specifications: a test on 22 agile product backlogs showed that LLMs outperformed standard NLP pipelines and managed to produce more correct diagrams, but also highlighted the importance of human validation, due to the high variability of the generated output. 

The study above, as well as most of the following ones, make use of PlantUML\footnote{\url{https://plantuml.com/}} to represent UML diagrams: the notation allows the definition of diagrams through a plain text description language where specific keywords represent UML elements and relationships, with interpreters parsing such textual content to produce a graphical representation of a diagram. The textual nature of PlantUML makes it particularly suited for LLM interaction, since language models produce text while struggling to generate images correctly (particularly during 2023-2024, when image generation capabilities were still in the early stage of development); a drawback of the notation, however, is that modifying diagrams requires one to change the textual descriptions and then render them again, making it slightly less user-friendly compared to free form modeling canvases.

Similar issues were encountered by Câmara et al. \cite{camara2023assessment}, whose experience report on ChatGPT's ability to generate UML class diagrams found that the model was able to produce syntactically correct diagrams with, however, relevant semantic violations such as incorrect association multiplicities and wrong relationship types; the inconsistent output across multiple queries with the same set of requirements was also mentioned as a cause of concern, and drove the definition of a conceptual framework for evaluating LLM-generated diagrams, stating that using LLMs for producing diagrams is feasible as long as prompts are designed carefully and humans validate the output. The experience report was followed by a second study \cite{camara2024benchmarks}, where the authors proposed a benchmarking framework for evaluating LLMs on software modeling activities to enable consistent and reproducible comparisons. More broadly, evaluating structured LLM output in cases where correctness is not limited to lexical and syntactical fidelity but must also consider semantic validity in relation to a domain model benefits from hallucination detection strategies \cite{huang2025hallucination}, according to which faithfulness measurement could be considered a valid proxy for the quality of structured output generation. These recent studies point to a convergence in the research community toward multi-dimensional evaluation criteria that assess syntax quality, semantic correctness, and adherence to domain constraints independently of one another.

A comparative study by Chen et al. \cite{chen2023automated} assessed GPT-3.5 against GPT-4 on ten modeling exercises, finding that LLMs were capable of generating relatively high quality classes and attributes but struggled the most with structural constraints related to associations (type, multiplicity, attribute placing). Li et al. \cite{li2024llm} investigated whether the use of Chain-of-Thought prompting could improve class diagram generation compared to simple LLM interaction with guided human extraction. Few-shot CoT prompting was the best method, outperforming guided human extraction in terms of class identification, and struggled the most with attribute identification, suggesting a division of concerns where LLMs propose the base structure and human experts provide refinement assistance. 

Other modeling notations were also covered beyond class diagrams: UML sequence diagram generation by ChatGPT was studied by Ferrari et al. \cite{ferrari2024model}, observing that the generation based on 28 sets of natural language requirements yielded satisfactory models in terms of lexical consistency and syntactic adherence to UML standards, while also pointing out that significant completeness and correctness problems arise when the requirements used to generate models are ambiguous or inconsistent. 

An assessment of four models (GPT-3.5, GPT-4, LLaMA-2, Mixtral) focused on generating Data Flow Diagrams from user stories \cite{herwanto2024automating} found that the proprietary GPT-4 achieved the best results in terms of completeness, while also pointing out the comparable performance of the open source option Mixtral, indicating that less general-purpose models may be an option for integrating LLMs in modeling related tasks, while reinforcing the importance of human oversight through dedicated editing interfaces.

Schleifer et al. \cite{schleifer2024minimal} defined two pipelines for generating use case diagrams from requirement specifications, observing the effects of named entity recognition combined with active learning against prompt engineering with generative models, finding that the generation process is feasible even with minimal training data and that LLMs can generalize over different application domains without requiring specific fine tuning.

Studies have also explored the use of LLMs for modeling tasks as assistive tools instead of autonomous model generators: an example is the study by Wang et al. \cite{wang2024llms}, who had 45 undergraduate students generate use case, class, and sequence diagrams with LLM support. Findings showed that LLM integration can accelerate the drafting process, while still requiring expert intervention and correction due to the presence of semantic and structural errors.

Other studies are relevant due to their focus on the educational aspects of LLM use for modeling tasks. For example, Kasneci et al. \cite{kasneci2023chatgpt} performed a survey on the opportunities and challenges of LLMs in education, identifying three relevant potential applications such as the production of personalized feedback, the ability to generate tailored example content, and the potential creation of exercises, while also pointing out risks related to academic integrity tied to students using LLMs to produce answers or to students relying excessively on AI-generated content.

Sarsa et al. \cite{sarsa2022automatic} found that OpenAI Codex, one of ChatGPT's precursor technologies, could generate programming assignments and code explanation with low effort, with most of the generated output requiring minimal teacher editing before being ready to use, suggesting potential LLM use for the generation of exercises. 

A factor common to most of the studies discussed in this section is the use of natural language descriptions as input for generating modeling artifacts, with the complexity of these descriptions affecting the quality of the output diagrams. An example of metric that can be used to measure the readability of natural language text is the Flesch-Kincaid readability score \cite{solnyshkina2017flesch}, a formula based on sentence length and syllable count, with lower scores indicating text that requires high cognitive effort to be understood, and higher scores signifying more accessible prose. This readability score can be used as a proxy to characterize requirements specification, and explain variance in the quality of LLM-generated models, since more complex requirements may be harder to parse for agents.

\section{Empirical Research Methods}
\label{sec:empirical_methods}

The empirical studies described in the following chapters share a common methodological foundation which is introduced here to provide the necessary context for the descriptions of the ensuing experimentations. For all studies, their goals were systematically defined according to the Goal-Question-Metric (GQM) approach \cite{basili1994gqm}, a structured template that is used to specify a measurement goal according to which are the objects under study, for what purpose the analysis is being conducted, what is the quality focus, who are the involved stakeholders with their viewpoint, and the context under which the study is executed. Following the definition of a study's goal, research questions are then derived from it, metrics are defined to answer these questions and, in most cases, null hypotheses are formulated around the goal to guide the statistical analyses that will answer the research questions.  Thanks to how it enforces traceability between research objectives and data collected, thus making the reasoning behind decisions explicit, the GQM template has become one of the standard framing devices for empirical software engineering research.

The controlled experiments described in Chapters 3 and 4 adopted within-subjects designs, studies in which all participants are exposed to more than a single experimental treatment, ensuring that each participant acts as their own control by observing all experimental conditions, in turn increasing the statistical power of the analysis. The two large-scale controlled studies employed two treatments and two tasks combined, leading to the use of a 2x2 full factorial crossover design \cite{vegas2016crossover}, which is used to balance the order of tasks and treatments according to different participant groups in order to reduce learning effects and carryover biases. Furthermore, an additional measure was undertaken to counterbalance task order effect by applying balanced Latin squares to ensure participants can experience every combination of conditions. 

For experiments where multiple statistical hypotheses were tested simultaneously on the same database, the risk of obtaining at least one false positive should be considered, and the design of the analysis should counteract it: the corrective action adopted to limit the rate of such Type I errors was the use of the Bonferroni correction, which adjusts the significance threshold by dividing the desired $\alpha$ level by the number of independent tests performed, ensuring that each individual test is held to a stricter standard \cite{abdi2007bonferroni}.

The threats to validity for all empirical studies described in subsequent chapters are characterized according to the taxonomy established by Wohlin et al. \cite{wohlin2012experimentation}, which distinguishes four categories: threats to internal validity concern factors that may affect the relationship between the treatment and the result and other factors that have not been considered during the analysis; threats to external validity concern whether the results can be generalized outside of the experiment's context; threats to construct validity relate to how appropriate the metrics selected for analysis actually represent the constructs the study is intended to observe; threats to conclusion validity impact the correctness of the statistical reasoning.

The controlled experiments involving students also measured the participants' opinions through questionnaires which also included open-ended questions: answers to said questions were analyzed following the open coding procedure described in the Straussian Grounded Theory \cite{strauss1998basics}. This approach consists of reading each textual answer and assigning it a label that summarizes its topic, selecting the label from a growing set of existing topics or adding a new one to the set if no suitable one exists yet; a second iteration is then performed to ensure answers are assigned a suitable topics, eventually assigning a new one out of the finalized set or having answers that suit more topic receive multiple labels. The analysis produces a ranked set of themes that summarizes the participants' opinions without relying on predefined categories, an approach that is suitable for an exploratory analysis of student opinions that are not known in advance.

Lastly, questionnaires also included validated psychometric instruments to measure student perception across perceived usability, technology acceptance, and perception of the gamified experience. Usability was measured through the System Usability Scale (SUS) \cite{brooke1996sus}, a set of ten questions that alternate between positively and negatively expressed statements answered on a Likert scale ranging from 1 to 5, with the resulting score being computed with Equation \ref{eq:sus_back}:
\begin{equation}
    S = 2.5 \cdot \left( \sum_{\substack{i=1,3,5,7,9}} (x_i - 1) + \sum_{\substack{i=2,4,6,8,10}} (5 - x_i) \right)
    \label{eq:sus_back}
\end{equation}

\noindent where $x_i$ is the raw score assigned to question $i$. Odd questions are those with a positive connotation, and their contribution is computed as $x_i - 1$, while even questions are negative ones, and their contribution is computed as $5 - x_i$. The total sum of the adjusted score is then multiplied by 2.5 to normalize the result to the $[0, 100]$ range. Conventionally, a score of 68 is considered the sign of average usability.

Technology acceptance was instead measured using the Technology Acceptance Model (TAM) \cite{davis1989perceived, lee2003technology}, a set of questions divided into four constructs that measure the Perceived Usefulness (PU), Perceived Ease of Use (PE), Behavioral Intention to Use (BI), and Attitude Towards Usage (ATU) of a system, with answers expressed as Likert-scale values ranging from 1 to 5. The gamification-specific perception was instead measured through the GAMEX questionnaire \cite{eppmann2018gameful}, an instrument designed specifically to measure the quality of a gamified environment beyond engagement counts through six constructs that measure aspects related to enjoyment, absorption, creative thinking, dominance, activation, and presence of possible negative effects.